\scanpage{29}
$=$ si $A$ courbe elliptique sur corps alg. clos ; donc
l'isom $\alpha$ n'est \uline{pas} fonctoriel p.r. à des
hom de schémas abéliens de $=$ dimension
[dans la mesure où ça pourrait avoir un
sens, les deux foncteurs n'étant pas de $=$
variance en $A$ \ldots]. Cependant, \deleted{\unsure{si $u$ est}
\unsure{un hom de $A,B$, on a}}
\[
\boxed{\operatorname{d\acute et} T_p(u)_{\mathrm{\acute et}}\,\operatorname{d\acute et} T_p(u^*)_{\mathrm{\acute et}}=\deg u}
\]
\margin{vaut aussi (via \unsure{identific.|identif.} qu'il permet) pour hom de Sch.A. de $=$ \unsure{dimension|dimensions} rel. $A,B$.}
% slanted

c) Il faudrait vérifier que
\[
\alpha_{A^*}=\alpha_A^{-1}
\]
[sans signe (?)\ldots]. Noter aussi le
comportemt par somme de VA. \ldots

d) L'isom $({}_pA)_{\mathrm{\acute et}}\mathbin{\to}({}_pA^*)_{\mathrm{\acute et}}$ induit par
$\alpha_A$ est bien celui déduit plus haut de
la dualité en cohomologie de DeRham
(à vérifier \ldots).
% archive sheet number: 28
